% ── Interlude 03: The Observational Frame ────────────────────
% Object class: window, threshold, boundary space
% Appears after Part III (Payne) before Part IV (future strand)
% No named characters. No date. No location.

\begin{interlude}{The Observational Frame}

There is an edge where the inside meets the outside.

This edge is not neutral. It selects. Of everything that exists
on the other side, only what can pass through the frame enters.
The rest remains outside, not absent, but unavailable---present
in the world and absent from the record simultaneously.

What the frame admits depends on the frame. This is obvious and
consistently overlooked. The observer positioned behind the frame
sees what comes through and calls it the world. The qualification
---that this is the world as the frame permits it to appear---is
noted, if at all, in a footnote, and then forgotten in the
interpretation.

The frame is not mentioned again.

It becomes the condition of seeing rather than one of its
components, transparent in the way that instruments are always
assumed to be transparent, contributing nothing to the image
while shaping everything about it.

\medskip

Every frame has a preferred orientation.

It admits more easily from some directions than others. Light at
certain angles passes through without loss. Light at other angles
is reflected, absorbed, scattered. The asymmetry is a physical
property of the frame's material, its thickness, the conditions
under which it was made.

What this means is that the world does not present itself equally
through the frame. Some things arrive clearly. Others arrive
degraded, stripped of certain frequencies, altered in ways that
may or may not be recoverable. Others do not arrive at all.

The observer does not experience this asymmetry directly.

The observer experiences only what arrives, and what arrives seems
complete---a full account of whatever is on the other side. The
gaps left by the frame's selectivity are not visible as gaps.
They are visible as absence, which is different. Absence suggests
that there is nothing there. A gap suggests that something is
there but was not transmitted.

Most frames produce gaps that look like absences.

\medskip

The frame also works in the other direction.

The observer is not only receiving. The act of observation changes
the conditions on the other side of the frame---subtly, sometimes,
in ways that compound over time and eventually become impossible
to ignore. The light sent through the frame to illuminate what is
being observed carries energy. The instrument positioned at the
boundary exerts pressure. The categories brought to the frame
determine what counts as signal and what counts as noise, and this
determination is not passive.

A frame is therefore not a window in the simple sense.

It is a site of exchange. Something passes inward; something passes
outward. The observer and the observed are in contact through the
frame, and both are altered by the contact, and the record produced
reflects this alteration without always recording it.

What is described as pure observation is always also intervention.

The frame is where this becomes unavoidable.

\medskip

There are frames that have been in place so long that they are no
longer recognized as frames.

They have become part of the architecture---built into the walls,
assumed in every account of what can be seen from this position.
To question them requires stepping back far enough to see the
frame itself, which means stepping out of the position the frame
defines, which means giving up the view it provides.

Most observers do not do this.

The view is too useful, too familiar, too deeply embedded in what
it means to see from here.

And so the frame persists, selecting, distorting, exchanging,
shaping the record of the world while remaining absent from it---
present as a condition, invisible as an object.

It stands at the boundary.

It has always stood there.

Everything that is known passed through it to become known.

\end{interlude}
